\documentclass[../../../include/open-logic-section]{subfiles}

\begin{document}

\olfileid{sth}{card-arithmetic}{ch}

\olsection{فرضیهٔ پیوستار}

نتیجهٔ پیشین به‌درستی نشان می‌دهد که توان‌رسانی کاردینالی بسیار
\emph{آسان} می‌بود اگر تضمین داشتیم که کاردینال‌های نامتناهی با توان‌های
$2$ «سرراست رفتار می‌کنند»، یعنی بنا بر
\olref[opps]{lem:SizePowerset2Exp} با گرفتن مجموعهٔ توانی سرراست رفتار
می‌کنند. اما نمی‌توانیم فرض کنیم که کاردینال‌های نامتناهی واقعاً با
مجموعه‌های توانی سرراست رفتار می‌کنند.

برای آغازِ واکاوی این نکته، نمادگذاری مناسبی معرفی می‌کنیم.

\begin{defn}
اگر $\cardsucc{\cardfont{a}}$ کوچک‌ترین کاردینالِ اکیداً بزرگ‌تر از
$\cardfont{a}$ باشد، دو دنبالهٔ نامتناهی زیر را تعریف می‌کنیم:
\begin{align*}
	\aleph_{0} &\defis \omega &
	\beth_{0} &\defis \omega\\
	\aleph_{\alpha \ordplus 1} &\defis \cardsucc{(\aleph_{\alpha})} &
	\beth_{\alpha+1} &\defis \cardexpo{2}{\beth_{\alpha}}\\
	\aleph_{\alpha} &\defis \bigcup_{\beta< \alpha} \aleph_{\beta} &
	\beth_{\alpha} &\defis \bigcup_{\beta < \alpha}\beth_{\beta} &
	\text{هرگاه $\alpha$ عددی ترتیبیِ حدی باشد.}
	\end{align*}
\end{defn}

تعریف $\cardsucc{\cardfont{a}}$ موجه است، زیرا
\olref[cardinals][classing]{lem:NoLargestCardinal} می‌گوید برای هر
کاردینال $\cardfont{a}$ کاردینالی بزرگ‌تر از آن وجود دارد، و استقرای
ترامتناهی تضمین می‌کند که یک \emph{کوچک‌ترین} کاردینالِ بزرگ‌تر از
$\cardfont{a}$ وجود دارد. تعریف‌های باقی‌ماندهٔ دنباله‌های $\aleph$ و
$\beth$ نیز به‌وسیلهٔ بازگشت ترامتناهی فراهم می‌شوند.

کانتور این نمادگذاری «$\aleph$» را معرفی کرد؛ این نماد \emph{الف}،
نخستین حرف الفبای عبری و نخستین حرف واژهٔ عبریِ «نامتناهی»، است. پیرس
نمادگذاری «$\beth$» را معرفی کرد؛ \emph{بِت} دومین حرف الفبای عبری
است.\footnote{پیرس این نمادگذاری را در نامه‌ای به کانتور در دسامبر
۱۹۰۰ به کار برد. متأسفانه، او در همان‌جا استدلالی نادرست نیز آورد مبنی
بر این‌که برای $\alpha \geq \omega$، $\beth_\alpha$ وجود ندارد.}
اکنون این نمادها کاردینال‌هایی نامتناهی در اختیار ما می‌گذارند.

\begin{prop}
برای هر عدد ترتیبی $\alpha$، $\aleph_\alpha$ و $\beth_\alpha$
کاردینال‌اند.
\end{prop}

\begin{proof}
هر دو نتیجه با یک استقرای ترامتناهی ساده به دست می‌آیند.
$\aleph_0 = \beth_0 = \omega$ بنا بر
\olref[cardinals][classing]{omegaisacardinal} یک کاردینال است. اگر
$\aleph_\alpha$ و $\beth_\alpha$ هر دو کاردینال باشند،
$\aleph_{\alpha+1}$ و $\beth_{\alpha+1}$ صریحاً به‌صورت کاردینال تعریف
شده‌اند. همچنین، بنا بر
\olref[cardinals][classing]{unioncardinalscardinal}، اجتماع مجموعه‌ای از
کاردینال‌ها خود یک کاردینال است.
\end{proof}
\noindent
افزون‌براین، هر کاردینال نامتناهی یک $\aleph$ است:

\begin{prop}
اگر $\cardfont{a}$ کاردینالی نامتناهی باشد، برای یک $\gamma$ یکتا داریم
$\cardfont{a} = \aleph_\gamma$.
\end{prop}

\begin{proof}
بر کاردینال‌ها استقرای ترامتناهی می‌کنیم. در گام استقرا فرض کنید هرگاه
$\cardfont{b} < \cardfont{a}$، داشته باشیم
$\cardfont{b} = \aleph_{\gamma_\cardfont{b}}$. اگر برای کاردینالی چون
$\cardfont{b}$ داشته باشیم $\cardfont{a} = \cardsucc{\cardfont{b}}$،
آنگاه $\cardfont{a} =
\cardsucc{(\aleph_{\gamma_\cardfont{b}})}=
\aleph_{\gamma_\cardfont{b}+1}$. اگر $\cardfont{a}$ جانشین هیچ
کاردینالی نباشد، چون کاردینال‌ها اعداد ترتیبی‌اند، داریم
$\cardfont{a} = \bigcup_{\cardfont{b} < \cardfont{a}} \cardfont{b} =
\bigcup_{\cardfont{b} < \cardfont{a}}{\aleph_{\gamma_\cardfont{b}}}$؛
پس $\cardfont{a} = \aleph_\gamma$، که در آن
$\gamma = \bigcup_{\cardfont{b} < \cardfont{a}}
\gamma_\cardfont{b}$. سرانجام، از تعریف و استقرای ترامتناهی نتیجه می‌شود
که نگاشت $\alpha\mapsto\aleph_\alpha$ اکیداً صعودی است؛ بنابراین چنین
$\gamma$ای یکتا است.
\end{proof}

از آن‌جا که هر کاردینال نامتناهی یک $\aleph$ است، این پرسش پیش می‌آید:
آیا هر کاردینال نامتناهی یک $\beth$ نیز هست؟ اگر \emph{چنین} می‌بود،
کاردینال‌های نامتناهی با عملِ گرفتن مجموعهٔ توانی «سرراست رفتار
می‌کردند». در واقع، آنگاه حکم زیر برقرار می‌بود:

\begin{defish}
\emph{فرضیهٔ تعمیم‌یافتهٔ پیوستار} (GCH). برای هر $\alpha$،
$\aleph_\alpha = \beth_\alpha$.
\end{defish}

افزون‌براین، اگر GCH برقرار بود، می‌توانستیم توان‌رسانی کاردینالی را
به‌طور چشمگیری ساده کنیم. به‌ویژه، می‌توانستیم نشان دهیم که هرگاه
$\cardfont{a}$ نامتناهی و $1 \leq \cardfont{b} < \cardfont{a}$ باشد،
مقدار $\cardexpo{\cardfont{a}}{\cardfont{b}}$ میان
$\cardfont{a}\leq \cardexpo{\cardfont{a}}{\cardfont{b}} \leq
\cardsucc{\cardfont{a}}$ محصور است. سپس می‌توانستیم شرایط دقیقی را بیان
کنیم که تعیین می‌کنند کدام‌یک از دو امکان رخ می‌دهد (یعنی آیا
$\cardfont{a} = \cardexpo{\cardfont{a}}{\cardfont{b}}$ یا
$\cardexpo{\cardfont{a}}{\cardfont{b}} =
\cardsucc{\cardfont{a}}$).\footnote{این شرط را مفهوم
\emph{هم‌پایانی} تعیین می‌کند.}

اما GCH یک \emph{فرضیه} است، نه یک \emph{قضیه}. در واقع،
\citet{Godel1938} اثبات کرد که اگر $\ZFC$ سازگار باشد، آنگاه
$\ZFC + \text{GCH}$ نیز سازگار است. ولی بعداً معلوم شد که به همان اندازه
می‌توان $\lnot$GCH را به $\ZFC$ افزود. در واقع، ساده‌ترین نمونهٔ
غیربدیهی GCH را در نظر بگیرید:

\begin{defish}
\emph{فرضیهٔ پیوستار} (CH). $\aleph_1 = \beth_1$.
\end{defish}

\citet{Cohen1963} اثبات کرد که اگر $\ZFC$ سازگار باشد، آنگاه
$\ZFC + \lnot\text{CH}$ نیز سازگار است. پس فرضیهٔ پیوستار از $\ZFC$
مستقل است.

نام «فرضیهٔ پیوستار» از آن‌جا می‌آید که «پیوستار» نام دیگری برای خط
اعداد حقیقی، یعنی $\Real$، است. از \olref[opps]{continuumis2aleph0}
داریم $\card{\Real} = \beth_1$. بنابراین فرضیهٔ پیوستار می‌گوید میان
کاردینالیتهٔ اعداد طبیعی، $\aleph_0 = \beth_0$، و کاردینالیتهٔ پیوستار،
$\beth_1$، هیچ کاردینالی وجود ندارد.

با توجه به \emph{استقلال} (G)CH از $\ZFC$، دربارهٔ \emph{صدق} آن‌ها چه
باید بگوییم؟ در این‌باره مطالب \emph{بسیاری} می‌توان گفت. در واقع، بخش
زیادی از پژوهش بارورِ اخیر در نظریهٔ مجموعه‌ها صرف بررسی همین مسائل شده
است. بااین‌حال، دست‌کم دو نکتهٔ بسیار کوتاه شایستهٔ تأکید است.

نخست: از این نتایج صوریِ استقلال، \emph{بی‌درنگ} نتیجه نمی‌شود که ارزش
صدق GCH یا CH \emph{نامتعین} است. چه‌بسا فقط لازم باشد اصول موضوعهٔ
بیشتری بیفزاییم که به نظرمان طبیعی‌اند و پرسش را به یکی از دو سو فیصله
می‌دهند. خود گودل نیز همین را پاسخ درست می‌دانست.

دوم: استقلال CH از $\ZFC$ بی‌شک \emph{شگفت‌انگیز} است، اما به معنای
لفظی کلمه \emph{باورنکردنی} نیست. نکته فقط این است که بنا بر آنچه
$\ZFC$ به ما می‌گوید، گذار از کاردینال‌ها به جانشین‌هایشان ممکن است به
ابزاری ظریف‌تر از صرفِ گرفتن مجموعه‌های توانی نیاز داشته باشد.
 %عمل گرفتن مجموعهٔ توانی ما را از یک مرحلهٔ سلسله‌مراتب به مرحلهٔ
 %جانشین آن می‌برد (از $V_{\alpha}$ به $V_{\alpha+1}$).

پس از این دو ملاحظه، اگر می‌خواهید بیشتر بدانید، باید به فیلسوفان و
ریاضی‌دانان گوناگونی رجوع کنید که در این مناقشه ذی‌نفع‌اند.\footnote{برای
آغاز شاید بد نباشد \citet[\S15.6]{Potter2004} را بخوانید.}

\end{document}
